\section{Start Here: Trainer and Server}

\begin{coachbox}{FOR THE TRAINER: Use the short path first}
In this guide, the trainer is the responsible adult who helps the server
practise. The priest or MC decides the liturgical route, and the adult
sacristan controls flame and hot charcoal.

First complete ``Today's Mass,'' let the server read the eight-scene map, and
turn to the assigned role sheet. Rehearse one scene at a time. Use the complete
team reference later only when the MC or trainer needs to check an exact cue,
handoff, or branch.

Before practice, an adult must demonstrate the actual church's steps, lanes,
object places, and reverences. Live flame, charcoal, the Communion plate, and
vessels require the supervision named in this guide.
\end{coachbox}

\begin{learnerbox}{Four things to know before you walk}
Write down your role. Point out where you start, the first object you touch,
where you wait during the Gospel, and the last object you put away. If you lose
your place, do not guess: stop, hold any object safely and still, watch the MC,
and rejoin at the next cue you know.
\end{learnerbox}

\begin{coachbox}{FOR THE TRAINER: Print one role packet}
Print this start sheet and the assigned role range below. Page numbers are the
numbers printed at the foot of the page; a shared boundary page is intentional.

\begin{center}
\small
\begin{tabularx}{\linewidth}{@{}>{\bfseries}l l >{\bfseries}l X@{}}
MC & \pageref{mc-role:mc:first}--\pageref{mc-role:mc:last} &
BB & \pageref{mc-role:bb:first}--\pageref{mc-role:bb:last}\\
A1 & \pageref{mc-role:a1:first}--\pageref{mc-role:a1:last} &
T1--T4 & \pageref{mc-role:t:first}--\pageref{mc-role:t:last}\\
A2 & \pageref{mc-role:a2:first}--\pageref{mc-role:a2:last} &
CB & \pageref{mc-role:cb:first}--\pageref{mc-role:cb:last}\\
TH & \pageref{mc-role:th:first}--\pageref{mc-role:th:last} &
L & \pageref{mc-role:l:first}--\pageref{mc-role:l:last}\\
\end{tabularx}
\end{center}

Keep the complete team reference for the MC or trainer. The response course is
on printed pages \pageref{course:first}--\pageref{course:last}. For detachable
practice, open the separate \emph{Missa Cantata Cue Cards} PDF from the same
catalog row. Its integrated response deck is on physical pages 1--8 and its
action deck is on pages 9--12; each complete range begins on an odd page.
\end{coachbox}

\subsection{Today's Mass}

Fill this in with the priest or MC before anyone vests. A blank is an unanswered
question, not permission to improvise.

\newcommand{\MCGuideCheck}{\fbox{\rule{0pt}{1.15ex}\rule{1.15ex}{0pt}}}
\newcommand{\MCGuideLine}[1]{\rule{#1}{0.4pt}}
\begin{center}
\small
\begin{tabularx}{\linewidth}{@{}>{\bfseries}p{0.29\linewidth}X@{}}
Server and assigned role & \MCGuideLine{0.61\linewidth}\\[5pt]
Route approved by (priest or MC) & \MCGuideLine{0.61\linewidth}\\[5pt]
Practice trainer & \MCGuideLine{0.61\linewidth}\\[5pt]
Incense & \MCGuideCheck\ yes: TH and BB assigned \quad
          \MCGuideCheck\ no: omit every incense row\\[5pt]
Entrance & \MCGuideCheck\ cross and CB \quad
           \MCGuideCheck\ candle pair without a cross\\[5pt]
Epistle & \MCGuideCheck\ priest sings \quad
          \MCGuideCheck\ assigned L sings from the prepared lectern\\[5pt]
Faithful Communion & \MCGuideCheck\ expected; plate prepared and A1 trained
                     \quad \MCGuideCheck\ not expected\\[5pt]
Torch condition & \MCGuideCheck\ no Communion and not a fast day: MC-041
                  \quad \MCGuideCheck\ Communion or fast day: MC-048\\[5pt]
Day's sung parts & Gloria: \MCGuideCheck\ yes \MCGuideCheck\ no \quad
                    Credo: \MCGuideCheck\ yes \MCGuideCheck\ no\\[5pt]
Ending / day branch & \MCGuideLine{0.61\linewidth}\\[5pt]
Named base, lanes, and object places & \MCGuideLine{0.61\linewidth}\\
\end{tabularx}
\end{center}

\subsection{The whole Mass in eight scenes}

Use this map to learn where you are in the Mass. Your role sheet tells you what
you do inside each scene; the action IDs lead your trainer to the complete team
reference when a handoff or branch needs more detail.

\newcommand{\MCSceneTile}[3]{%
  \begin{minipage}[t]{0.48\linewidth}
  \begin{tcolorbox}[
    enhanced,
    colback=black!3,
    colframe=black!38,
    boxrule=0.45pt,
    arc=0pt,
    left=6pt,
    right=6pt,
    top=4pt,
    bottom=4pt,
    height=0.145\textheight,
    valign=top]
    \textbf{\large #1}\enspace\textbf{#2}\par\smallskip
    {\small #3}
  \end{tcolorbox}
  \end{minipage}}

\noindent
\MCSceneTile{1}{Get ready}{Check the roster and objects, line up, process, land at
the named places, and account for the priest's biretta (MC-001--007).}\hfill
\MCSceneTile{2}{Begin}{The quiet response team serves the prayers at the foot; the
team then follows the Introit, Kyrie, Gloria when appointed, and Collect
posture (MC-008--018).}

\noindent
\MCSceneTile{3}{Listen}{Follow the Epistle route, prepare and serve the Gospel, sit
for the sermon if one is given, and take the appointed Credo gestures
(MC-019--029).}\hfill
\MCSceneTile{4}{Prepare}{Serve wine and water, the selected incense route and
Lavabo, and the quiet \textit{Suscipiat} (MC-030--035).}

\noindent
\MCSceneTile{5}{Adore}{Follow the Preface, bring torches, kneel for the Canon, and
use the exact bell and elevation-incense cues (MC-036--042).}\hfill
\MCSceneTile{6}{Receive}{Follow the Lord's Prayer and Agnus Dei; when the faithful
receive, use the bell and plate route; then complete both ablutions and the
book transfer (MC-043--050).}

\noindent
\MCSceneTile{7}{Finish}{Take the day's Postcommunion posture and dismissal
branch; add the blessing when it belongs and the Last-Gospel route when it is said
(MC-051--054).}\hfill
\MCSceneTile{8}{Return}{Collect only the named exit objects, make the assigned
final reverence, recess, and account for every person and object in the
sacristy (MC-055--057).}

\MCEntranceHandoffsDiagram

\section{Your Missa Cantata team}

A \textit{Missa Cantata} is a sung Mass without deacon and subdeacon. The
priest sings his appointed parts, the schola or choir sings its parts, and the
servers carry out the sanctuary work. That means there are two different kinds
of answering:

\begin{itemize}
  \item the designated servers make the quiet minister's answers at the foot of
  the altar and during the prayers the priest does not sing; and
  \item the choir and people make the public sung responses. A server joins
  those only when the director has assigned him to do so.
\end{itemize}

\begin{watchbox}{Three source labels}
\textsc{Missal} marks a rule stated in the 1962 Missal. \textsc{Model} marks
the ordinary ceremonial pattern used by this guide, based on the contemporary
Fortescue--O'Connell tradition. \textsc{Local choice} marks something the
priest or master of ceremonies must settle before rehearsal. Never turn a local
choice into a surprise during Mass.
\end{watchbox}

\subsection{The model roster}

\begingroup
\renewcommand{\rolecard}[2]{\item[\textbf{#1}] #2}
\begin{description}[leftmargin=0pt,labelindent=0pt,labelwidth=0pt,labelsep=0pt,style=nextline,
  topsep=0.2em,itemsep=0.25em]
\rolecard{MC --- master of ceremonies}{Knows the complete route, cues the
team quietly, assists the priest when assigned, and fixes problems without
drawing attention to them. In this model MC handles the biretta, serves the
Missal, and transfers the book after the ablutions. The MC is not a second
celebrant and does not make large theatrical signals.}
\rolecard{A1 and A2 --- acolytes}{Work as a matched pair. They carry the
processional or Gospel candles in this model, serve wine and water, serve the
Lavabo, and return every object to its prepared place. A1 is the model bell
and Communion-plate server; A2 brings the Lavabo water and basin.}
\rolecard{TH and BB --- thurifer and boat bearer}{Used only on the
with-incense track. TH keeps the thurible safe and ready; BB presents the boat.
This guide's incense route requires both named servers; it does not include a
combined-role shortcut.}
\rolecard{T1--T4 --- torchbearers}{These four servers are part of this guide's
core model. They enter empty-handed, receive four lighted torches at the end of
the Preface, kneel symmetrically for the sacred action, return the torches in
pairs, and recess empty-handed. A team using a different number needs a newly
marked roster and route.}
\rolecard{CB --- cross-bearer, when used}{Carries the processional cross
upright and steady. Some ordinary entrances have no cross-bearer; this is a
\LocalChoice.}
\rolecard{L --- assigned lesson singer, alternate branch}{The priest sings the
Epistle in the default model. If another qualified minister is assigned, L
uses MC-020 and is added to the roster before vesting.}
\end{description}
\endgroup

\begin{watchbox}{Why ``Today's Mass'' must be completed}
General Rubric 426 permits incense at a sung Mass but does not require it at a
Missa Cantata; this model assigns both TH and BB when it is used. The model's
quiet response team is MC, A1, and A2. A different assignment, the alternate L
route, a processional cross, Communion of the faithful, and every church-specific
place must be written and rehearsed before vesting. A1 handles the Communion
plate only after the required hands-on training.
\end{watchbox}

\section{Move like one team}

\subsection{The six basic actions}

\begin{description}[style=multiline,leftmargin=0.25\linewidth,labelwidth=0.22\linewidth]
  \item[Hands joined.] When empty-handed, hold the palms together at the chest,
  fingers straight and pointing upward. Do not swing the arms.
  \item[Head bow.] Bow the head smoothly at the name of Jesus and at the other
  names or moments taught by the MC. Return without a snap.
  \item[Body bow.] Bend from the waist, pause, and rise together. Use it only
  where this guide or the MC assigns it.
  \item[Genuflection.] Stop, keep the body upright, touch the right knee to the
  floor, rise, and then move. Object-bearers use the reverence assigned to that
  object; they do not improvise a genuflection while carrying a tall cross,
  lighted candle, or full thurible.
  \item[Turn and walk.] Turn toward the direction of travel. Walk forward,
  at a prayerful pace. In the entrance and recession, travel single file
  through the nave; partner positions apply only after the group forms at a
  station. Never walk backward down steps.
  \item[Give and receive.] Present an object at a usable height with its handle
  ready. Look at the receiver's hands. Release only when the receiver has it.
\end{description}

\begin{watchbox}{Reverence depends on what is present}
The tabernacle's location, whether the Blessed Sacrament is reserved or exposed,
and whether a server is carrying an object can change the proper reverence.
The MC must demonstrate the rule in the actual sanctuary. This guide does not
replace that site lesson. Lay servers do not add hand or object kisses unless
the director expressly teaches the local practice.
\end{watchbox}

\begin{watchbox}{When sung words require a genuflection}
If the proper text marks sung words for a genuflection, the MC gives the cue
and every free-handed server faces the altar and genuflects with the priest.
At the Gospel, TH also faces the altar and genuflects while holding the closed
thurible; A1 and A2 remain upright while holding their candlesticks.
Never guess from a familiar word: follow the proper rubric and the MC.
\MissalRule: General Rubrics 519 and 521(d).
\end{watchbox}

\subsection{Incense safety}

TH checks the chains, lid, charcoal, spoon, and route before vesting. Hold the
thurible away from albs, hair, woodwork, and other people. Never run with it,
leave it unattended on a combustible surface, or add charcoal without the
adult sacristan's permission. The MC teaches the number and direction of swings;
the young server's first job is safe, steady presentation.

\section{MC / Trainer Reference: Every Action in Order}

The following chart preserves the complete model route. A server may use an
action ID to check one detail, but the chart is chiefly for the MC and trainer,
who must coordinate every person, object, branch, and waiting state.

\begingroup
% This local treatment leaves the shared action-table contract untouched. The
% stronger stage banners and separate finish line make the long chronology
% usable as a dispatch sheet instead of a wall of equally weighted prose.
\renewcommand{\actionsection}[1]{%
  \needspace{7\baselineskip}%
  \par\medskip
  \noindent\colorbox{black!78}{%
    \parbox{\dimexpr\linewidth-2\fboxsep\relax}{%
      \color{white}\bfseries\strut #1}}\par\smallskip}
\renewenvironment{actiontable}{%
  \footnotesize
  \setlength{\LTpre}{0.08em}
  \setlength{\LTpost}{0.45em}
  \setlength{\tabcolsep}{5pt}
  \renewcommand{\arraystretch}{1.13}
  \begin{longtable}{@{}>{\RaggedRight\arraybackslash}p{0.18\linewidth}>{\RaggedRight\arraybackslash}p{0.14\linewidth}>{\RaggedRight\arraybackslash}p{0.62\linewidth}@{}}
  \toprule
  \textbf{Cue / action ID} & \textbf{Who moves} & \textbf{Action and finish state}\\
  \midrule
  \endfirsthead
  \toprule
  \textbf{Cue / action ID} & \textbf{Who moves} & \textbf{Action and finish state}\\
  \midrule
  \endhead
  \bottomrule
  \endfoot
}{%
  \end{longtable}
}
\renewcommand{\actionrowid}[5]{%
  \rule{0pt}{2.7ex}\textbf{\textsc{#1}}\par\smallskip #2 &
  \textbf{#3} &
  #4\par\smallskip\textbf{Finish / next:} #5\\[0.55em]}

\actionsection{Before the entrance}

\begin{actiontable}
\actionrowid{MC-001}{At the credence}{MC}{Check the Missal on its stand, wine
and water cruets, Lavabo basin and towel, bell, the Communion plate when the
faithful will receive, and the marked stands for two Gospel candles and four
torches. On the alternate lesson branch, place the marked Epistle book on its
lectern and assign the adult who will recover it after Mass. Confirm that the
priest or sacristan has prepared the chalice on the
altar; no young server carries it in this model. Check the biretta place at the
sedilia and name each object's pickup and return place.}{Nothing is merely
somewhere nearby.}
\actionrowid{MC-002}{With-incense track}{TH / BB}{Light charcoal early enough for a
steady coal. Test the lid and rehearse the route to the priest.}{Adult
sacristan approves the coal.}
\actionrowid{MC-003}{All vested}{MC}{Name every role and ask each server for his first
three jobs. Check that cassocks, surplices, and shoes are neat. Have each
server repeat the missed-cue rule: stop guessing, hold any object safely and
still, watch the MC, and rejoin at the next known cue.}{No whispered role
changes in the procession, and no invented recovery route.}
\actionrowid{MC-004}{Objects and line-up}{Adult sacristan; all lay servers}{The
adult lights the two processional candles and gives one to A1 and one to A2;
the pair checks equal height. If CB is used, CB takes the cross from its named
stand. Then take this single-file model order: TH, BB on the incense branch;
CB, A1, A2 when the cross is used, or A1, A2 when it is not; T1, T2, T3, T4;
L on the alternate lesson branch; MC; priest. Omitted roles close the gap.
Partners form side by side only at their stationary stations.}{The adult owns
every flame handoff. \LocalChoice: the church doorway, nave route, and
stationary formation can change spacing.}
\end{actiontable}

\actionsection{Entrance and prayers at the foot}

\begin{actiontable}
\actionrowid{MC-005}{Sacristy signal}{All lay servers}{Bow to the cross with the priest as directed,
turn together, pass through the doorway one behind another, and begin the
procession through the nave without a gap.}{Hands joined unless carrying an
object; travel single file.}
\actionrowid{MC-006}{At the sanctuary}{Object-bearers}{Stop on the marked
arrival line. At the MC's signal make the reverence assigned to the carried
object; then A1/A2 set candles at the credence, CB sets the cross in its stand,
and TH/BB go to the safe incense station. T1--T4 and L go to their base
places.}{Do not bunch at the center; every object now has a place.}
\actionrowid{MC-007}{At the sanctuary}{MC / free-handed servers}{When the
priest removes the biretta, MC receives it with both hands. Form the rehearsed
straight line and make the prescribed reverence with the priest. MC sets the
biretta at the sedilia and meets the priest at the center foot; every other
free-handed server clears the middle and returns to base.}{\ModelRule: no lay
server uses ceremonial kisses; the marked arrival line is \LocalChoice.}
\actionrowid{MC-008}{Foot prayers begin}{MC / A1 / A2}{Kneel in the prepared
places and make the quiet answers together, at the priest's pace. Sign
yourself at \textit{In nomine Patris}; bow the head at \textit{Gloria Patri};
then sign yourself again at \textit{Adiutorium
nostrum}. Stay upright during the priest's \textit{Confiteor}; turn and bow
toward him for \textit{Misereatur tui}. Bow for your \textit{Confiteor}, turn
toward him at \textit{tibi, pater} and \textit{te, pater}, and strike your
breast three times at \textit{mea culpa}. Stay bowed through his
\textit{Misereatur vestri}; rise upright and sign yourself during
\textit{Indulgentiam}; then bow slightly through the final versicles.}{Use the
response course for the words. Do not compete with the choir's Introit.}
\actionrowid{MC-009}{At a Mass of the season when Psalm 42 is omitted}{MC / A1 / A2}{Keep
\textit{Introibo} and its answer, then continue at \textit{Adiutorium}.}{\MissalRule:
where this ordinary route is otherwise applicable, this occurs from Passion
Sunday through Holy Thursday. Requiem Mass is excluded from this guide.}
\actionrowid{MC-010}{Priest ascends}{MC}{Rise at the rehearsed cue, clear the
center, and take the marked book-side place. A1/A2, T1--T4, and L remain at
base; TH/BB remain at the safe station until the incense cue.}{A step is never
crossed blindly.}
\end{actiontable}

\actionsection{Introit through the Collect}

\begin{actiontable}
\actionrowid{MC-011}{After the altar kiss, with incense}{TH / BB / MC}{Approach
at the Epistle side. BB gives the boat to the priest and receives it back; TH
gives him the open thurible. MC keeps the route clear. After the priest has
incensed the altar, MC receives the thurible securely while TH waits at MC's
left.}{\MissalRule: the priest follows the Solemn-Mass incensation.}
\actionrowid{MC-012}{First incensation ends}{MC / TH / BB}{MC and TH make the
rehearsed bows; MC incenses the priest with three double swings and gives the
thurible to TH. TH closes it and returns with BB and the boat to the safe
station.}{\MissalRule: at Missa Cantata the priest is incensed. The giver,
receiver, count, and return are this guide's \ModelRule.}
\actionrowid{MC-013}{No-incense track}{All lay servers}{Omit every incense row. Take the places
assigned for the Introit and Kyrie.}{Do not imitate an empty hand-off.}
\actionrowid{MC-014}{Kyrie is sung}{MC}{Stand at the book-side place. If the
priest goes to the sedilia, escort him by the marked lane, give him the
biretta with both hands after he sits, and stand beside the seat. Near the
final invocation, receive the biretta, replace it at the sedilia, and return
with him. A1/A2 remain at base.}{\MissalRule:
the priest may sit while the Kyrie is sung; the lane is \LocalChoice.}
\actionrowid{MC-015}{Gloria intonation}{All lay servers}{Stand with hands
joined when empty. Bow the head at \textit{Deo} in the intonation and, during
the sung text, at \textit{Adoramus te}, \textit{Gratias agimus tibi}, \textit{Iesu Christe},
and \textit{Suscipe deprecationem nostram}; sign yourself at the final
\textit{Cum Sancto Spiritu}. Join the public chant only if assigned.}{Skip
the entire Gloria route when it is not appointed.}
\actionrowid{MC-016}{During the sung Gloria}{MC}{If the priest sits, use the
complete MC-014 escort-and-biretta sequence, remain in attendance, and return
with him for the final phrase.}{Skip the whole row when there is no Gloria.}
\actionrowid{MC-017}{\textit{Dominus vobiscum}}{Choir / people; all lay
servers}{The public sung answer is \textit{Et cum spiritu tuo}. All servers
stand at base, hands joined when empty; they do not add a second answer unless
assigned to the public chant.}{Sung response, not the quiet Low-Mass voice.}
\actionrowid{MC-018}{Collects}{Choir / people; all lay servers}{The public
answer to each appointed conclusion is \textit{Amen}; servers sing only if
assigned. All stand for the sung greeting. On Sundays and feasts remain
standing, hands joined when empty, through the Collects. At ferial Masses of
Advent, Lent, and Passiontide, September Ember days, and second- or third-class
vigils outside Paschaltide, kneel after the greeting and rise after the final
Collect.}{\MissalRule: General Rubric 521 gives the named kneeling branch.
More than one Collect occurs only as the day's rubrics allow.}
\end{actiontable}

\actionsection{Epistle, Gospel, and Creed}

\begin{actiontable}
\actionrowid{MC-019}{Epistle --- priest sings}{All lay servers}{Stand at base,
hands joined when empty, and listen. Do not add the Low-Mass \textit{Deo
gratias} after the chanted lesson.}{This is the guide's default lesson route.}
\actionrowid{MC-020}{Epistle --- L sings}{L / MC}{At the MC's cue, L goes
empty-handed from base to the prepared lectern, makes the site's prescribed
reverence on the route, sings the Epistle, makes the same reverence, and
returns to base. MC cues but does not shadow him; the marked book remains on
the lectern until the named adult recovers it after the concluding sacristy
reverence at MC-057.}{\MissalRule: a minister may sing the Epistle. This alternate must
be selected before vesting.}
\actionrowid{MC-021}{Gradual / Alleluia / Tract}{MC}{If the priest sits, MC
escorts him by the MC-014 lane and returns with him at the agreed cue before
the Gospel preparation. Every other server stands at base, hands joined when
empty.}{Do not guess from the length of the chant.}
\actionrowid{MC-022}{Gospel preparation, with incense}{TH / BB / MC}{At the
MC's cue, TH and BB come from the safe station to the Epistle side. BB gives
the boat to the priest and receives it back; TH gives him the open thurible and
receives it after he has put in and blessed incense. BB returns the boat to the
safe station. TH and MC go to the center for the Gospel movement.}{The complete
out-and-back handoff is part of the incense branch.}
\actionrowid{MC-023}{Gospel candles}{A1 / A2}{Take the two lighted candles from
their credence stands together, go to the center, make the prescribed
object-bearer's reverence, turn together, and go level to the marked Gospel-side
places.}{\ModelRule: the candle pair flanks the Gospel-side altar book.}
\actionrowid{MC-024}{Gospel movement}{MC / TH / A1 / A2}{The Missal remains on
its Gospel-side altar stand. Form without crossing paths: TH ahead if used,
A1 and A2 level with their candles, and MC beside the book. Face the book and
hold still.}{No incense: omit TH; the candle pair and MC keep the same route.}
\actionrowid{MC-025}{Gospel title and text}{Choir / people; all lay servers}{The
public answer is \textit{Gloria tibi, Domine}. Free-handed servers make the
three small signs of the Cross; candle-bearers keep the candles upright. On the
incense track, TH gives the thurible to MC, MC gives it to the priest, then MC
receives it and returns it to TH after the book is incensed. TH holds it closed
at the marked place. If a Gospel word requires a genuflection, MC, TH, and
every other free-handed server face the altar and genuflect with the priest;
A1 and A2 remain upright with the candlesticks.}{\MissalRule: General Rubric
519 supplies the positive genuflection and the candle-acolyte exception.
Servers sing only if assigned to the public response.}
\actionrowid{MC-026}{Gospel ends}{TH / A1 / A2 / MC}{Do not add the Low-Mass
\textit{Laus tibi, Christe} after the chanted Gospel. A1 and A2 return the
candles to their credence stands together. On the incense track TH returns the
closed thurible to its safe station; MC clears the route and all return to
base.}{The priest is not incensed after he sings the Gospel.}
\actionrowid{MC-027}{Sermon, if given}{All lay servers}{At the MC's cue, go
empty-handed to the named server bench and sit. At the closing cue, stand
together and return to base, leaving the center clear.}{Rejoin MC-028 when a
Credo follows; otherwise rejoin MC-030.}
\actionrowid{MC-028}{Credo and \textit{Et incarnatus est}}{All lay servers}{Join the chant only if
assigned. Bow the head at \textit{Iesum Christum} and \textit{simul
adoratur}. When \textit{Et incarnatus est} is sung, follow the rehearsed branch: if the
priest is standing, genuflect with him; if he is seated, make a profound head
bow. At the three Christmas Masses and the Annunciation, all kneel.}{\MissalRule:
General Rubrics 518--519. MC pre-cues the branch. Free-handed servers sign
themselves at \textit{Et vitam venturi saeculi}.}
\actionrowid{MC-029}{During the Credo}{MC}{If the priest sits, use the full
MC-014 escort-and-biretta sequence and remain beside the sedilia. At the final
phrase, return with him. Every other server remains at base and follows
MC-028.}{Skip
if there is no Credo.}
\end{actiontable}

\MCGospelFormationDiagram

\actionsection{Offertory}

\begin{actiontable}
\actionrowid{MC-030}{Offertory begins}{A1 / A2}{Go together from base to the
credence. A1 takes the wine cruet; A2 takes the water cruet. With handles ready
for the receiver, go side by side to the marked Epistle-side handoff point.}{Do
not reverse cruets or sides.}
\actionrowid{MC-031}{Wine and water}{A1 / A2}{A1 gives wine to the priest and
receives it back; then A2 gives water and receives it back. Make the prescribed
reverence together and return both cruets to their marked credence places.}{The
same giver receives the same cruet.}
\actionrowid{MC-032}{Offertory incense begins}{TH / BB / MC; all other lay
servers}{TH and BB approach
the Epistle side. BB gives the boat to the priest and receives it back; TH
gives the open thurible. After the gifts and altar are incensed, MC receives
the thurible, incenses the priest with three double swings, and gives it to TH.
As soon as the priest has been incensed, A1 and A2 begin MC-033 on the inner
credence lane while TH uses the outer lane. This core model has no clergy
seated in choir; if clergy are present, the MC must supply their rank and
count on a separate sheet before rehearsal. After A1 and A2 return, TH uses
one double swing for each recipient or pair in this complete model roster, in
order: MC; A1 and A2 as a pair; BB; T1 with T2; T3 with T4; then optional CB
and L separately. Each recipient or pair faces TH and bows before and after. TH then
incenses the people with three double swings---center, Gospel side, Epistle
side---closes the thurible, and returns with BB to the safe station.}{\MissalRule:
incense is permitted; the priest, ministers, and people are incensed. The
complete no-clergy roster order, counts, pairs, and noncollision lanes are this
guide's declared \ModelRule.}
\actionrowid{MC-033}{Priest reaches the Lavabo place}{A1 / A2}{On the incense
track, begin as soon as the priest has been incensed, while TH continues the
remaining MC-032 incensations. On the no-incense track, begin at MC's cue when
the priest comes to the Epistle-side Lavabo place after preparing the chalice.
A2 holds the basin and pours
water over the priest's fingers; A1 presents and receives the towel. Bow
together and return the set by the inner credence lane before TH reaches the
server-recipient line. If delayed, hold clear and let TH wait; never cross his
outer route.}{\MissalRule: the priest washes after the offering and, when
incense is used, after he has been incensed. The two cues are explicit.}
\actionrowid{MC-034}{\textit{Orate, fratres}}{MC / A1 / A2}{Wait until the
priest has finished the invitation and turned back, then say
\textit{Suscipiat Dominus} together and clearly.}{This is not sung.}
\actionrowid{MC-035}{Secret conclusion}{Choir / people; all lay servers}{The
public answer is \textit{Amen}. All servers stand at base, hands joined when
empty; they sing only if assigned.}{Responses are in the course below.}
\end{actiontable}

\actionsection{Preface, Canon, and elevations}

\MCElevationTableauDiagram

\begin{actiontable}
\actionrowid{MC-036}{Preface dialogue}{Choir / people; all lay servers}{The
public responses are \textit{Et cum spiritu tuo}, \textit{Habemus ad Dominum},
and \textit{Dignum et iustum est}. All servers stand; at \textit{Gratias agamus
Domino Deo nostro}, free-handed servers bow the head. Sing only if assigned and
use the chant tone, not the sound-spelling rhythm.}{T1--T4 remain at base until
the end-of-Preface cue.}
\actionrowid{MC-037}{End of Preface}{MC / adult sacristan / T1--T4}{The adult
sacristan lights four torches at their marked stands. At MC's cue, T1--T4 each
take the assigned torch, form two pairs, go to the center by the marked lane,
make the prescribed object-bearer's reverence together, divide evenly, and
kneel at their four marked places.}{\ModelRule: four torches and this route;
live-flame work remains under adult supervision.}
\actionrowid{MC-038}{Sanctus and Canon}{All lay servers; A1}{At the opening
\textit{Sanctus}, A1 gives this model's one short bell signal, sets the bell on
its marked pad, and kneels at the assigned step. All other non-torch servers
kneel at their marked places; T1--T4 are already kneeling with torches.}{\MissalRule:
the small bell is rung at the \textit{Sanctus}; one short signal and the exact
stations are the declared \ModelRule.}
\actionrowid{MC-039}{Before Consecration and at elevations}{A1}{Shortly before
the Consecration, give one warning signal at the MC's visual cue. At the Host
elevation ring three distinct strokes; repeat three at the chalice elevation.
Keep the bell low and still between signals and replace it on its pad after the
second group.}{\MissalRule: a warning signal is given shortly before the
Consecration, and each elevation has three strokes or a continuous ring. This
guide selects three distinct strokes.}
\actionrowid{MC-040}{Before and at each elevation, with incense}{TH / BB / MC}{At
MC's cue before the Consecration, TH opens the thurible and BB gives the boat
and spoon to MC. MC puts in incense without a blessing and gives the boat back
to BB. BB returns it to the safe station; TH takes the marked Epistle-side
place and kneels. TH incenses the Host three times at its elevation and the
chalice three times at its elevation, then closes the thurible and returns it
to the safe station.}{\MissalRule supplies the unblessed incense and three
incensations at each elevation; the exact handoff and station are this guide's
\ModelRule.}
\actionrowid{MC-041}{Each elevation}{All lay servers}{Look toward the altar prayerfully. Bow when
the MC cues, then return upright together. After the chalice elevation and the
priest's final genuflection, the ordinary no-Communion torch branch rises at
MC's cue and completes its paired return to the four stands. Each torch is set
securely, the adult extinguishes it, and T1--T4 return empty-handed to base;
the extended Communion or fast-day branch remains kneeling.}{The ordinary
departure happens here; the extended branch finishes at MC-048. No extra
signals or movements.}
\actionrowid{MC-042}{Canon doxology}{Choir / people; all lay servers}{The
public answer is \textit{Amen}. After that answer, as the Lord's Prayer begins,
non-torch servers rise and stand at base, hands joined when empty. Only T1--T4
on the extended Communion or fast-day branch remain kneeling with torches
until MC-048; the ordinary branch is already empty-handed at base.}{Servers sing only
if assigned. The Canon posture begun at MC-038 has now reached the
\textit{Pater noster}.}
\end{actiontable}

\actionsection{Communion}

\MCAblutionsDiagram

\begin{actiontable}
\actionrowid{MC-043}{Lord's Prayer}{Choir / people; all lay servers}{The public
response is \textit{Sed libera nos a malo}; the later conclusion receives
\textit{Amen}. Non-torch servers stand at base, hands joined when empty;
extended T1--T4 remain kneeling and hold the torches upright; the ordinary
torch branch stands empty-handed at base.}{Do not turn the full
\textit{Pater noster} into a server response.}
\actionrowid{MC-044}{\textit{Pax Domini}}{Choir / people; all lay servers}{The
public answer is \textit{Et cum spiritu tuo}. Servers keep the MC-043 posture
and do not add another reply unless assigned to the chant.}{This core model
does not add a server pax route.}
\actionrowid{MC-045}{Agnus Dei and priest's Communion}{All lay servers}{During
the three sung invocations, each free-handed server makes one breast strike;
torchbearers still holding torches keep both hands on them. At the MC's cue all non-torch
servers kneel. Make the two rehearsed body bows at the priest's reception of
the Host and Precious Blood; object-bearers make only the safe inclination
taught for their object.}{Join the chant only if assigned. Requiem is outside
this ordinary route.}
\actionrowid{MC-046}{Before faithful Communion}{A1; all other lay servers}{If
the faithful will receive, A1 gives one short warning signal at the MC's cue,
sets the bell on its pad, and then prepares for the plate pickup. Every other
server kneels; T1--T4 keep the torches. There is no second \textit{Confiteor},
\textit{Misereatur}, or \textit{Indulgentiam}: the printed rite goes directly
to \textit{Ecce Agnus Dei}.}{\MissalRule: \textit{Ritus servandus} X.6. Skip
the whole row when no one else receives.}
\actionrowid{MC-047}{Faithful Communion}{A1}{At the MC's cue, take the
Communion plate from its prepared place. Accompany the priest on the rehearsed
side, keep the plate level under each communicant's chin, and watch the Host rather
than the line. After the last communicant, carry it level to the altar and give
it to the priest for the required fragment check and transfer; receive the
cleared plate only when he releases it, then return it upright to its marked
credence place. Alert the priest at once if anything falls.}{\MissalRule: the
plate is used beneath communicants and any fragments are placed in the chalice.
The priest's side and approach lane are \LocalChoice; this job requires
hands-on adult training.}
\actionrowid{MC-048}{First ablution and extended torch return}{A1 / A2;
T1--T4 / adult sacristan}{A1 and A2 bring both cruets to the Epistle-side
handoff point. A1 pours the first wine into the chalice. The ordinary
no-Communion torch branch already returned at MC-041. If the faithful received,
or on a fast day, T1--T4 have remained through Communion; after this first
ablution, at the MC cue they rise, turn, make the prescribed object-bearer's
reverence, and return in pairs to the four named torch stands. Set each torch
securely; the adult sacristan extinguishes it. The second ablution may continue
without waiting for their sacristy round trip.}{Do not leave one torch alone.
The fast-day extension and first-ablution departure apply the Solemn-Mass
baseline as this guide's \ModelRule.}
\actionrowid{MC-049}{Second ablution}{A1 / A2}{Remain at the handoff point after
MC-048. A1 pours wine and A2 pours water over the priest's fingers into the
chalice. Receive the cruets, bow together, and return both to their marked
credence places.}{Move quietly while the Communion chant continues.}
\actionrowid{MC-050}{Book transfer}{MC}{After the ablutions, take the Missal
and its stand from the Gospel side of the altar, check that the center is
clear, carry both with two hands across the front route, make the prescribed
reverence, and set them on the Epistle side at the Introit angle. Return to the
book-side place.}{\MissalRule: the minister transfers the book after the
ablutions. This model binds the task to MC.}
\end{actiontable}

% This last chronology heading needs room for its repeated table head and the
% opening row; the ordinary action-section reserve is too short at this exact
% page boundary.
\actionsection{Postcommunion, Last Gospel, and recession}

\begin{actiontable}
\actionrowid{MC-051}{Postcommunion and Prayer over the People}{Choir / people;
A1 / A2; all other lay servers}{The public responses are \textit{Et cum
spiritu tuo} and \textit{Amen}; servers sing only if assigned. All stand for
the sung greeting. On Sundays and feasts remain standing, hands joined, for
the Postcommunion. At ferial Masses of Advent, Lent, and Passiontide,
September Ember days, and second- or third-class vigils outside Paschaltide,
kneel after the greeting for the Postcommunion. On a ferial Mass from Ash
Wednesday through Wednesday of Holy Week, when this ordinary route is
otherwise applicable, remain kneeling while the priest sings \textit{Humiliate
capita vestra Deo} and the Prayer over the People; answer its conclusion only
if assigned, then rise for the final greeting. A1 and A2 confirm that exit
objects remain stable before returning to base.}{\MissalRule: General Rubric
521 and \emph{Ritus servandus} XI.2 govern the named posture and Prayer-over-
the-People branch. Do not pack anything still needed.}
\actionrowid{MC-052}{Dismissal}{Choir / people; all lay servers}{For the normal
\textit{Ite, missa est}, the public response is \textit{Deo gratias}; at a
Mass of the season within the Easter octave, both add \textit{alleluia} twice.
When a procession is to
follow, \textit{Benedicamus Domino} receives \textit{Deo gratias}; skip
MC-053--054 and use the separate procession sheet. Servers stand at base and
sing only if assigned.}{Requiem is outside this ordinary guide.}
\actionrowid{MC-053}{Blessing}{Choir / people; all lay servers}{On the normal
ending, all servers kneel together at MC's cue and sign themselves as the
priest blesses. The public answer is \textit{Amen}; servers sing it only if
assigned. Rise together afterward.}{\MissalRule: \textit{Benedicamus Domino}
has no blessing.}
\actionrowid{MC-054}{Last Gospel}{MC / A1 / A2; all other lay servers}{On the
normal route, MC, A1, and A2 make the quiet answers \textit{Et cum spiritu
tuo}, \textit{Gloria tibi, Domine}, and finally \textit{Deo gratias}. All
free-handed servers make the three small signs of the Cross at the title and
genuflect with the priest at \textit{Et Verbum caro factum est}. At the third
Mass of Christmas, the Last Gospel is omitted: go directly to MC-055.}{Other
proper or omitted Last-Gospel cases belong to their edition-checked special
sheet.}
\actionrowid{MC-055}{Form the recession}{MC / TH / BB / CB / A1 / A2 /
T1--T4 / L}{After the Last Gospel, or at the omission cue, MC checks the center
and takes the priest's biretta from its marked sedilia place, holding it flat
with both hands. On the incense branch TH takes the closed thurible and BB the
boat; CB takes the cross when used; A1 and A2 take the two candles; T1--T4 and
L remain empty-handed with hands joined. Form the same single-file order used
at MC-004; use the locally marked paired formation only for the stationary
reverence.}{The priest's biretta and every entrance object have now rejoined the
exit chain. The adult sacristan has confirmed whether charcoal remains live.}
\actionrowid{MC-056}{Final reverence and recession}{All lay servers}{At MC's
signal, free-handed T1--T4 and L genuflect if the Blessed Sacrament is reserved
at this altar and make a profound bow if it is not. This model keeps MC, TH,
BB, CB, A1, and A2 upright with their objects and assigns them a head bow. MC
then gives the biretta to the priest with both hands. Turn in the rehearsed
direction without crossing, then recess in the same order and spacing as the
entrance, one behind another through the nave and sacristy doorway.}{\MissalRule: \emph{Ritus servandus} XII.6 supplies the reserved/
unreserved branch. \ModelRule: the exact object-safe bows and biretta handoff.}
\actionrowid{MC-057}{In the sacristy}{All lay servers / adult sacristan}{Form
the same straight line and make the customary concluding reverence with the
priest. TH and BB give the thurible and boat to the adult at the metal-safe
station; A1 and A2 set candles in their stands; CB sets the cross in its stand;
T1--T4 help only with the objects named before Mass. On the alternate lesson
branch the assigned adult now recovers the Epistle book from its lectern and
returns it to its named sacristy place while L reports that it is accounted
for. The adult handles live charcoal and extinguishes flames. Report any
damage or mistake honestly.}{Serving
ends when every person and object is accounted for.}
\end{actiontable}
\endgroup

\section{Branch card: know what changes}

\begin{coachbox}{Ordinary branches covered here}
\begin{description}[style=multiline,leftmargin=0.27\linewidth,labelwidth=0.24\linewidth]
  \item[No incense.] Omit TH, BB, every incensation, and every incense hand-off.
  \item[No Gloria/Credo.] Skip the corresponding intonation, seating, and
  return rows.
  \item[Psalm 42 omitted.] Keep the opening \textit{Introibo} exchange and
  continue at \textit{Adiutorium}; do not omit the whole foot-prayer course
  unless General Rubric 424 applies.
  \item[Requiem.] This form is outside the guide. Psalm 42, the Agnus Dei text,
  dismissal, blessing, and some ceremonial actions differ. Rehearse from an
  edition-checked Requiem sheet.
  \item[Easter dismissal.] At a Mass of the season within the Easter octave,
  the prescribed dismissal and response add \textit{alleluia} twice.
\end{description}
\end{coachbox}

\begin{watchbox}{Special rites need their own sheet}
The Asperges, Candlemas, Ash Wednesday, Palm Sunday, Holy Week, Rogation
processions, Nuptial Mass, funerals and absolution, Mass before the Blessed
Sacrament exposed, pontifical functions, and Communion outside the ordinary
route are not compressed into this chart. The MC must supply the proper
edition-checked supplement.
\end{watchbox}
